\chapter{Virtual Reality}

Virtual Reality (VR) is a new interface paradigm (not an application) that allows users to interact with a computer or data based on an spatial interface, realtime rendering and realtime multi-modal interaction. There is a continuous 3D representation, and the holograhic space is scaled $1:1$. Two key aspects of VR are immersion and presence.
\begin{itemize}
    \item Immersion: Objective description of what any particular system does provide. It is the extent to which the VR system can deliver an inclusive, extensive, surrounding and vivid illusion of reality to the senses of a human participant. It is a property of the system, not of the user.
    \item Presence: is a subjective state of consciousness, the (psychological) sense of being in the virtual environment. It is a property of the user, not of the system. It is the extent to which the user feels present in the virtual environment, rather than in the physical environment.
\end{itemize}

Virtual Reality is a new interface paradigm, and as such, it has its own terminology. The most important concepts are:
\begin{itemize}
    \item \ul{VR System}: is the system that implements the VR paradigm. It is made up of hardware and software components.
    \item \ul{Virtual Object}: is a 3D computer generated object that the user can interact with.
    \item \ul{Real World}: the physical world that we live in.
    \item \ul{Virtual World}: a 3D computer generated scene made up of one or more virtual objects as well as their behavior and interaction between the objects. 3D user interfaces and avatars are also part of the virtual world.
    \item \ul{Virtual Environment}: is the combination between the VR system and the virtual world.
\end{itemize}

Once we have defined the key concepts of VR, the four basic elements of VR are:
\begin{itemize}
    \item Virtual World
    \item Immersion
    \item Interactivity
    \item Realtime feedback from the VR-System
\end{itemize}

In any of these elements is not present, it is hardly going to be a VR system. Some examples are:
\begin{itemize}
    \item Computer game: they have a virtual world, realtime interaction, realistic effects, realtime rendering... but there is no immersion, and it is not a specialized VR system. Therefore, it is not a VR system.
    \item Online Worlds: they may be realistic with immense functionality and visual effects... but no immersion is provided. Therefore, they are not VR systems.
\end{itemize}
    
There are $3$ different types of VR:
\begin{itemize}
    \item \ul{Virtual Reality (VR)}: is the paradigm where everything that the user sees is virtual.
    \item \ul{Augmented Reality (AR)}: is the paradigm where the user sees the real world, but with virtual objects superimposed on it.
    \item \ul{Mixed Reality (MR)}: it is a paradigm where the user sees a combination of the real world and virtual objects, but the virtual objects are anchored to the real world and can interact with it.
\end{itemize}

In order to achieve a good VR experience and a seamless integration of the user into the virtual world, there are several requirements that must be met:
\begin{itemize}
    \item Synchronization of many different components, such as sensors, trackers, rendering engines, display devices, etc. 
    \item Fast response of the system to user inputs, a crucial aspect for maintaining immersion and presence.
    \item Fast output of internal states of the system.
\end{itemize}

Obviously, VR is a realtime system that strives to minimize latencies. The latency of a VR system can be defined as the time it takes for the system to respond to a user input and update the virtual world accordingly. It inclues the user's reaction time, the sensor measurement time, the time it takes to update the internal state of the system, the rendering time and the display device specific latencies.\\

VR has some properties, such as dispkaying in scale $1:1$ or a fast and intuitive navigation, that provides several different applications, as mistakes identification (which is usually easier in 3D), safe simulation of dangerous situations (which allows military, aviation or space training) or even as an ``Eye Catcher'' at trade fares.

\section{History of VR}

Some important milestones in the history of VR are:
\begin{itemize}
   \item 1966: The first head-mounted display (HMD) was created by Ivan Sutherland, called the Sword of Damocles. It was a very primitive device that was suspended from the ceiling and had a very limited field of view.
   \item 1983: The first data glove was created. It was a glove that had sensors that could detect the movement of the fingers and the hand, allowing for more natural interaction with virtual objects. 
   \item 1992: CAVE (Cave Automatic Virtual Environment) was created. It is a room-sized VR system that uses projectors to display the virtual world on the walls, floor and ceiling, allowing for a more immersive experience.
   \item 2003: The first virtual driving simulator was created. It was a VR system that allowed users to experience driving a car in a virtual environment, providing a safe and controlled environment for training and testing.
\end{itemize}

However, an important date in the history of VR is 2016, when VR is popularized and a lot of companies start to invest in VR technology, such as Oculus, HTC, Sony, etc. This is due to the fact that the technology has advanced enough to provide a good VR experience at a reasonable cost.

\section{Motion Capturing}

Tt is the process of recording the movement of objects or people. It has evolved during time, and these were some important techniques used in the past:
\begin{itemize}
    \item Chronophotography: multiple photographs of a moving object were taken at regular intervals, allowing to analyze the motion of the object.
    \item 2D Rotoscoping: a technique where firstly a live-action film is shot, and then the frames are traced over to create an animation. It was used in the early days of animation, as it allowed for more realistic movement of characters. Some important examples of this technique are the movies Snow White and the Seven Dwarfs (1937), Pinocchio (1940) or Peter Pan (1953).
    \item 3D ``handmade'' mocap (motion capture): a technique where the movement of a character is created by hand, using keyframes to define the movement of the character. However, the movements are not painted by hand, but a software creates the movement of the character based on the keyframes. The person only decides how is the skeleton position at a given time, and the software creates the movement of the character based on the keyframes.
\end{itemize}

Nowadays, however, the motion capture is done using markers at each joint of the body, and the movement of the markers is recorded and processed to create the movement of the character. This technique is called marker-based motion capture, and it is widely used in the film and video game industry. There are different types of markers, which are usually attached to the body of the actor and may even be mixed:
\begin{itemize}
    \item Acoustic: they emit sound waves that are picked up by microphones. Depending on the time it takes for the sound waves to reach the microphones, the position of the markers can be calculated.
    \item Inertial: they use accelerometers and gyroscopes to measure the movement of the markers.
    \item LED: they emit light that is picked up by cameras.
    \item Magnetic: they use magnetic fields to measure the position of the markers.
    \item Reflective: they reflect light that is emitted by cameras. 
\end{itemize}



\section{Tracking}

In order to use those markers introduced in the previous section, we need to track their position in space. Tracking is the process of determining the position and orientation of objects in space. It is a crucial aspect of VR, as it allows for the correct registration between real and virtual objects. Two steps are necessary for tracking:
\begin{itemize}
    \item Calibration: the offline determination of all invariant parameters, such as the position of fixed sensors or the fixed camera properties.
    \item Tracking: the online determination of all time-varying parameters, such as the position of people or real objects.
\end{itemize}

During the following sections, we will see different tracking techniques.

\subsection{Acoustic Tracking}

In 2D, two circles can uniquely
determine two points, so a third circle is needed to uniquely determine just one pont. In 3D, two spheres can uniquely determine a whole circle, and a third sphere can uniquely determine two points.\\

In order to determine the position of a point $P$ in 3D space, let's firstly suppose that we have three fixed points $P_1$, $P_2$ and $P_3$ (the centers of the spheres), and we know the distance $r_i$ from each of these points to $P$ (the radius of the spheres). The position of $P$ can be determined by finding the intersection of the three spheres.
\begin{itemize}
    \item Let's first do a naïve approach, and do not do any transformation of the coordinate system. Then, $P=(x,y,z)$, and each of the centers of the spheres is given by $P_i=(x_i,y_i,z_i)$. The system of equations that we need to solve is shown in Equation~\eqref{eq:acoustic_tracking_naive}.
    \begin{equation}
        \begin{cases}
            (x-x_1)^2+(y-y_1)^2+(z-z_1)^2=r_1^2\\
            (x-x_2)^2+(y-y_2)^2+(z-z_2)^2=r_2^2\\
            (x-x_3)^2+(y-y_3)^2+(z-z_3)^2=r_3^2
        \end{cases}
        \label{eq:acoustic_tracking_naive}
    \end{equation}

    As we can see, this system of equations is non-linear, and it is not easy to solve.

    \item Therefore, the coordinate system is carefully chosen. The center of the first sphere is placed at the origin, so $P_1=(0,0,0)$. The $x$-axis is defined by the line that goes from $P_1$ to $P_2$, so $P_2=(d,0,0)$, where $d$ is the distance between $P_1$ and $P_2$. The $y$-axis is perpendicular to the $x$-axis and contained in the plane defined by $P_1$, $P_2$ and $P_3$, so $P_3=(i,j,0)$. Therefore, the $z$-axis is also perpendicular to the $x$-axis and $y$-axis, and it is defined by the right-hand rule. With this coordinate system, the system of equations that we need to solve is shown in Equation~\eqref{eq:acoustic_tracking}.
    \begin{equation}
        \begin{cases}
            x^2+y^2+z^2=r_1^2\\
            (x-d)^2+y^2+z^2=r_2^2\\
            (x-i)^2+(y-j)^2+z^2=r_3^2
        \end{cases}
        \label{eq:acoustic_tracking}
    \end{equation}

    Taking the first equation minus the second equation, we get:
    \begin{equation*}
        (x^2+y^2+z^2)-[(x-d)^2+y^2+z^2]=r_1^2-r_2^2
        \Longrightarrow
        2dx - d^2=r_1^2-r_2^2
        \Longrightarrow
        x=\frac{r_1^2-r_2^2+d^2}{2d}
    \end{equation*}

    Taking the first equation minus the third equation, we get:
    \begin{multline*}
        (x^2+y^2+z^2)-[(x-i)^2+(y-j)^2+z^2]=r_1^2-r_3^2
        \Longrightarrow
        2ix+2jy-i^2-j^2=r_1^2-r_3^2
        \Longrightarrow \\ \Longrightarrow
        y=\frac{r_1^2-r_3^2+i^2+j^2-2ix}{2j}
    \end{multline*}

    Finally, we can substitute the values of $x$ and $y$ in the first equation to get the value of $z$.
    \begin{equation*}
        z=\pm\sqrt{r_1^2-x^2-y^2}
    \end{equation*}

    Therefore, the position of $P=(x,y,z)$ in that given coordinate system is given by:
    \begin{equation*}
        \begin{cases}
            x=\dfrac{r_1^2-r_2^2+d^2}{2d}\\ \\
            y=\dfrac{r_1^2-r_3^2+i^2+j^2-2ix}{2j}\\ \\
            z=\pm\sqrt{r_1^2-x^2-y^2}
        \end{cases}
    \end{equation*}
    As we can see, two possible positions of $P$ are obtained.

    \item A last simplification can be done. Given that the points $P_1$, $P_2$ and $P_3$ can be placed in any position in space, we can carefully place them in such a way that their coordinates are very simple: $P_1=(0,0,0)$, $P_2=(1,0,0)$ and $P_3=(0,1,0)$. With this coordinate system, the system of equations that we need to solve is also very simple, and the position of $P$ is given by:
    \begin{equation*}
        \begin{cases}
            x=\dfrac{r_1^2-r_2^2+1}{2}\\ \\
            y=\dfrac{r_1^2-r_3^2+1}{2}\\ \\
            z=\pm\sqrt{r_1^2-x^2-y^2}
        \end{cases}
    \end{equation*}
\end{itemize}

\subsubsection{Determining the distance $r_i$}

Therefore, our only remaining problem is to determine the distance $r_i$ from each of the fixed points to $P$. It is calculated by measuring the time it takes for a sound wave to travel from the fixed point to $P$ (or vice versa). Given that the speed of sound is known, the distance can be calculated as $r_i=v\cdot t$, where $v$ is the speed of sound and $t$ is the time it takes for the sound wave to travel from the fixed point to $P$. However, some problems that may arise are:
\begin{itemize}
    \item The speed of sound is not constant, as it depends on the temperature, humidity and pressure of the air. Therefore, it is necessary to measure these parameters in order to calculate the speed of sound accurately.
    \item It requires a clear line of sight between the fixed points and $P$, as the sound waves can be reflected or absorbed by objects in the environment, which can lead to inaccurate measurements.
    \item Requires an additional trigger to decide when to start measuring the time. This trigger may be a button press, a light sensor, etc. However, this trigger can introduce additional latencies in the system.
    \item Multiple trackers can interfere with each other, as the sound waves emitted by one tracker can be picked up by the microphones of another tracker, leading to inaccurate measurements. Therefore, it is necessary to use different frequencies for each tracker, or to use a round-robin approach, where each tracker is activated one at a time.
\end{itemize}

An additional approach that can be used to determine the distance $r_i$ is to use the phase difference of the sound waves emitted and received. However, a $\phi$ phase can't be distinguished from a $\phi+2\pi$ phase, so it is not exactly accurate. A $\unit[40]{KHz}$ sound wave has a wavelength of $\unit[8.275]{mm}$, so it helps up to this accuracy.\\

This technique is widely used in VR systems, as it is simple to evaluate and the trackers can be small and lightweight. However, it depends on the environment and gives only position, not orientation (which would represent the normal vector).

\subsection{Inertial Tracking}

Inertial tracking is a technique that uses accelerometers and gyroscopes to measure the movement of the markers. It is based on the First Law of Newton, the Law of Inertia, which states that an object at rest will stay at rest, and an object in motion will stay in motion with the same speed and in the same direction unless acted upon by an unbalanced force.

\subsubsection{Acceleration sensor}

Let's first suppose that we can measure the aceleration of the marker, $a(t)$. We can therefore calculate the position of the marker as follows:
\begin{equation*}
    p(t)=p(0)+\int_0^t v(\tau)d\tau=p(0)+\int_0^t\int_0^\tau a(\sigma)d\sigma d\tau
\end{equation*}
where we have assumed that the initial velocity of the marker is $0$. Therefore, the position of the marker can be calculated by integrating the acceleration twice. In order to calculate the acceleration, the Second Law of Newton can be used, which states that the force $F$ acting on an object is equal to the mass $m$ of the object multiplied by its acceleration $a$, so $F=ma$. An object with a constant mass $m$ is attached to a spring, and the force $F$ is measured by the deformation of the spring. Therefore, the acceleration can be calculated as $a=\nicefrac{F}{m}$.
\begin{itemize}
    \item In reality, an acceleration sensor is an integrated circuit built from weigths and piezo elements, which are elements that generate an electric charge when they are deformed. The acceleration of the marker causes the deformation of the piezo elements because of the inertia of the weights, and the electric charge generated by the piezo elements is proportional to the acceleration of the marker. Therefore, the acceleration can be calculated by measuring the electric charge generated by the piezo elements.
\end{itemize}

\subsubsection{Gyroscope}

A gyroscope is a device that measures the orientation of an object in space. It is based on the principle of conservation of angular momentum, which states that an object will maintain its orientation in space unless acted upon by an external torque. A gyroscope consists of a spinning rotor (which is a heavy disk) that is mounted on a gimbal (which is a set of rings that allow the rotor to rotate freely in any direction). Even if the base of the gyroscope is moved, the rotor will maintain its orientation in space due to the conservation of angular momentum. Therefore, by measuring the orientation of the rotor, the orientation of the object to which the gyroscope is attached can be determined.
\begin{itemize}
    \item In reality, a gyroscope is an integrated circuit built from a vibrating structure. When the orientation of the gyroscope changes, the Coriolis force causes the vibrating structure to deform, and the electric charge generated by the deformation is proportional to the angular velocity of the gyroscope. Therefore, the orientation of the object can be calculated by measuring the electric charge generated by the deformation of the vibrating structure.
\end{itemize}~\\

Inertial tracking is widely used in VR systems, as they are self contained small and cheap devices, and can be used simultaneously. Devices usually have the IMU (Inertial Measurement Unit), which is a combination of an accelerometer and a gyroscope, and it can provide both position and orientation of the marker. However, a limitation is that each sensor can only measure one DOF (Degree of Freedom). Acceleration has $3$ DOF (the three components of the acceleration vector), and the gyroscope has $3$ DOF (the three components of the angular velocity vector). Therefore, 6 sensors are needed to measure the 6 DOF of the marker (3 for position and 3 for orientation).
\subsection{Magnetic Tracking}

Magnetic tracking is a technique that uses magnetic fields to measure the position of the markers. Some electrical coils are placed in the environment, and they generate a magnetic field. The markers have sensors that can measure the intensity of the magnetic field, and by measuring the intensity of the magnetic field at different points in space, the position of the marker can be determined taking into account that the intensity of the magnetic field decreases cubic with distance. 

Sensors are really small, do not require a line of sight and have very high update rates. However, they are very sensitive to metal objects in the environment, cables are required and have limited range.


\subsection{Optical Tracking}

Optical tracking is a technique that uses cameras to measure the position of the markers. The markers can be active (they emit light) or passive (they reflect light). The cameras capture the light emitted or reflected by the markers, and by analyzing the images captured by the cameras, the position of the markers can be determined. There are different types of optical tracking techniques, such as:
\begin{itemize}
    \item \ul{Single camera vs. Multi camera}: a single camera can only capture the position of the markers in 2D, while multiple cameras can capture the position of the markers in 3D by triangulation.
    \item \ul{Marker-based vs. Markerless}: marker-based tracking uses markers that are attached to the objects or people being tracked, while markerless tracking uses computer vision techniques to track the objects or people without the need for markers.
    
    Markless optical tracking requires advanced computer vision techniques, such as SIFT (Scale-Invariant Feature Transform). SIFT is a technique that detects and describes local features in images. Multiple images of the same scene but with different gaussian blur and different scales are taken, and the local features are detected in each image. The local features are then described by a vector of 128 dimensions, which is invariant to scale and rotation. Therefore, by comparing the local features detected in different images, the position of the objects or people can be determined.
\end{itemize}


Two main techniques are widely used in optical tracking: outside-in and inside-out.
\subsubsection{Outside-in}
In outside-in tracking, there are fixed decives (usually cameras) in the environment that track the position of the markers. They are oriented towards the user, and they capture the light emitted or reflected by the markers. Two main devices that use this technique are:
\begin{itemize}
    \item \ul{Oculus Rift}: it uses a single camera to track the position of the markers, which are active (they emit light). The camera captures the light emitted by the markers, and by analyzing the images captured by the camera, the position of the markers can be determined. Usually googles and controllers are worn by the user.
    \item \ul{HTC Vive - Lighthouse}: multiple laser emitters are placed in the environment, and they emit laser beams that sweep across the environment. The markers have sensors that can detect when they are hit by the laser beams, and by measuring the time it takes for the laser beams to hit the markers, the position of the markers can be determined.
\end{itemize}

The main advantages of outside-in tracking are that it has high precision and unlimitated theoretical range. However, it requires multiple lines of sight between the cameras and the markers, and it is computationally expensive, as it requires processing the images captured by the cameras to determine the position of the markers.

\subsubsection{Inside-out}

In inside-out tracking, the cameras are attached to the user pointing towards the environment, recognizing the environment and tracking the position of the markers. The main device that uses this technique is Oculus Quest 2.